\section{Related Work}
\label{sec:related}

\subsection{Graph-based Anomaly Detection}
Graph neural networks have been studied to detect anomalies in various domains, such as detecting review spams in business websites~\cite{dou2020enhancing,li2019spam,liu2020alleviating, liu2021pick}, 
rumors in social media~\cite{bian2020rumor,wu2020graph,yang2020rumor}, fake news~\cite{AA-HGNN2020,FangModel}, financial fraud~\cite{liu2019geniepath,xu2021towards,rao2021xfraud}, insurance fraud~\cite{liang2019uncovering}, and bot fraud~\cite{yao2020botspot}.
Most of them target how to design a proper aggregator that can distinguish the effects of different neighbors and reduce the inconsistency issue~\cite{liu2020alleviating} during message passing. For example, GAS~\cite{li2019spam} adopts the vallina GCN~\cite{kipf2016semi}. SemiGNN~\cite{SemiGNN} and Player2Vec~\cite{zhang2019key} apply attention mechanisms to assign low weights to suspicious links. To thoroughly reduce the negative influence of the suspicious links, several attempts~\cite{franceschi2019learning, liu2020alleviating} have been made to remove the suspicious links before graph convolution. CARE-GNN~\cite{dou2020enhancing} further adopts reinforcement learning to sample links according to the suspicious likelihood.
None of them are aware of the limitations caused by the objective function. To our knowledge, we are the first to change the commonly-used binary classification into the graph contrastive learning paradigm. 

\subsection{Graph Pre-training Schemes}
With the advances of self-supervised (pre-training) techniques in visual representation learning~\cite{chen2020simple,he2020momentum, chen2021exploring}, graph pre-training schemes have also attracted increasing attention. A naive GNN pre-training scheme is to reconstruct the vertex adjacency matrix, and GAE~\cite{kipf2016variational} and GraphSAGE~\cite{Graphsage} are two representative models. In addition to preserve the structure homophily, GPT-GNN~\cite{hu2020gpt} preserves the attribute homophily by predicting the masked node attributes. Motivated by~\cite{hjelm2018learning}, DGI~\cite{velickovic2019deep} and Infograph~\cite{Infograph} have been proposed to maximize the mutual information between the embeddings of the entire graph and the node within it. 
GraphCL~\cite{you2020graph} maximizes the mutual information between the embeddings of two graph instances, which are obtained from the same graph via graph data augmentation. Later, GCC~\cite{qiu2020gcc} shrinks the contrast between graphs into that between ego-networks, where the ego-network instances are obtained via random walk with start from the concerned node. All of them are proven to be useful for the downstream node classification task, but are not specifically proposed to solve the anomaly detection problem. DCI~\cite{wang2021decoupling} is a cluster-based version of DGI, 
which maximizes the fine-grained mutual information between the embeddings of a cluster and the nodes within it. DCI is proposed for graph anomaly detection. However, it is thoroughly unsupervised model and different from the corrupted graphs by \pmodel, the anomalies in DCI are thoroughly disconnected from the normal nodes, 
which may reduce the learning difficulty of the pseudo labels.

